\subsection{Task 1, Question 1}
\label{sec:q1}
\begin{quote}
For the audio encoder case, how many bytes does the kernel read and write in total? Which fraction of your GPU's peak bandwidth did you reach?
\end{quote}
The harness counts bytes moved as
$2 \cdot (\text{q\_rows} + \text{k\_rows}) \cdot \text{rotary\_dim} \cdot 4$
(read $+$ write over the rotated prefix). For the audio-encoder case
($B{=}1$, $\mathrm{seq}{=}1500$, $H_q{=}H_k{=}20$, $\text{rotary\_dim}{=}32$),
$\text{q\_rows}{=}\text{k\_rows}{=}30{,}000$, so
$2 \cdot 60{,}000 \cdot 32 \cdot 4 = 15.36 \times 10^6$~bytes (${\sim}14.6$~MiB)
per call (an analytical count, not a new timing measurement).
The fresh harness reports 0.026~ms and 600.0~GB/s, \textbf{78\%} of the
quoted 768~GB/s peak---a roofline-style~\cite{williams2009roofline} check
consistent with a memory-bound kernel that also executes trigonometry and indexing.

\subsection{Task 2, Question 3}
\label{sec:q2}
\begin{quote}
How many times does your softmax read each element of \texttt{x} from global memory? How many times does the textbook version (max, then sum, then write) read it?
\end{quote}
Our softmax reads each element \textbf{twice}: the first loop in
\texttt{softmax\_kernel} merges \texttt{MaxSum} statistics; the final output loop
rereads \texttt{xr[c]} to compute its normalised exponential.
The textbook implementation reads it \textbf{three times}, once each for maximum,
exponential sum, and output. Online softmax merges the first two passes, but
does not remove the output pass. These are source-level global-load requests,
not measured DRAM transactions; cache hits can reduce actual DRAM traffic.

\subsection{Task 3, Question 2}
\label{sec:q3}
\begin{quote}
Your kernel reaches only a fraction of cuBLAS's speed on the large cases. Name one thing cuBLAS does differently that explains part of the gap. You do not have to implement it.
\end{quote}
A plausible explanation is that cuBLAS likely uses more sophisticated
hierarchical tiling and register-level reuse, computing multiple outputs per
thread rather than our single-output $16\times16$ shared-memory tile.
Such strategies increase reuse and arithmetic intensity; however, cuBLAS is
closed-source and we did not inspect the particular kernel dispatched here.
Our encoder \texttt{q\_proj} reaches 2967~GFLOP/s versus cuBLAS's 20305;
\texttt{fc1 + GELU} reaches 2870 versus 21346~GFLOP/s.
The comparison is not a pure epilogue match (our \texttt{fc1} fuses GELU),
but the large library gap is consistent with the profile's dominant linear cost.
We do not claim tensor-core use by the particular cuBLAS calls without tracing them.

\subsection{Task 4, Question 1}
\label{sec:q4}
\begin{quote}
How much memory would the full score matrix of one encoder layer take (20 heads, 1500 x 1500, float32)? How much shared memory does one block of your kernel use instead?
\end{quote}
The full score matrix takes $20\cdot1500\cdot1500\cdot4
=180{,}000{,}000$ bytes ($171.66$~MiB).
In \texttt{launch\_attention}, shared memory is allocated as
$4[32(d+1)+32d+4d]$ bytes: padded keys, values and cached scaled queries.
For encoder head dimension $d=64$, this is
$4[32\cdot65+32\cdot64+4\cdot64]=17{,}536$ bytes
($17.125$~KiB) per block, \emph{not} 48~KiB.
The latter is only the launcher's threshold for requesting an opt-in shared-memory
limit. Scores are reduced tile by tile with online softmax rather than storing the
full matrix; these memory footprints are analytical counts from the source.
